\documentclass[11pt,a4paper]{article}
\usepackage[utf8]{inputenc}
\usepackage[T1]{fontenc}
\usepackage{geometry}
\geometry{margin=1in}
\usepackage[hidelinks]{hyperref}
\usepackage{enumitem}
\usepackage{titlesec}
\usepackage{xcolor}
\usepackage{listings}
\usepackage{fancyhdr}
\usepackage{parskip}

\pagestyle{fancy}
\fancyhf{}
\rhead{CodeTrail User Manual}
\lhead{Version 1.0}
\cfoot{\thepage}

\title{\textbf{CodeTrail User Manual}}
\author{CodeTrail Team}
\date{2026-11-16}

\begin{document}

\maketitle

\begin{center}
\textbf{Product:} CodeTrail — Collaborative Coding Platform \\
\textbf{Document:} User Manual \\
\textbf{Version:} 1.0 \\
\textbf{Date:} 2026-11-16
\end{center}

\tableofcontents
\newpage

% ==========================================
\section{Introduction}

CodeTrail is a collaborative coding platform that combines a modular workspace, real-time collaboration, sandboxed code execution, and contribution telemetry/audit features in one web application.

\textbf{Key capabilities}
\begin{itemize}
  \item \textbf{Authentication:} Email/password and Google-based sign-in (as configured in your deployment).
  \item \textbf{Home Dashboard:} Personalized overview of recent activity, workspaces, metrics, and goals.
  \item \textbf{Workspaces:} Create, join, search, filter, star, and manage collaborative coding rooms.
  \item \textbf{Modular Workspace:} Browser-based editor with file management, collaboration, and execution.
  \item \textbf{Code Execution:} Run code in an isolated sandbox and view stdout, stderr, compilation errors, exit codes, and execution time.
  \item \textbf{Telemetry \& Integrity Audit:} Activity tracking with proof-of-work/chain verification views.
  \item \textbf{Analytics:} Contribution and session insights (where available).
\end{itemize}

% ==========================================
\section{System Requirements \& Access}

\begin{itemize}
  \item \textbf{Supported browsers:} Latest versions of Google Chrome, Microsoft Edge, or Mozilla Firefox (recommended: Chromium-based browser).
  \item \textbf{Internet connection:} Required for authentication, collaboration, execution, and data sync.
  \item \textbf{Screen size:} Desktop/laptop recommended for the full workspace/editor experience; responsive layouts are provided where supported.
  \item \textbf{Access URL:} Use the CodeTrail URL provided by your deployment/administrator.
\end{itemize}

\begin{quote}
\textbf{Note:} Some features (execution, AI assistance, email, OAuth) depend on backend services and environment configuration. If a service is unavailable, the UI may show an offline/error state.
\end{quote}

% ==========================================
\section{Getting Started}

\subsection{Landing Page}
When not logged in, CodeTrail shows the public landing page with product overview, feature highlights, and options to \textbf{Log In} or \textbf{Sign Up}.

\subsection{Sign Up}
\begin{enumerate}
  \item From the landing/login screen, choose \textbf{Sign Up} (or switch to the sign-up form).
  \item Enter the required details (e.g., name, email, password) as prompted.
  \item Submit the form to create your account.
  \item After successful registration, sign in to access your dashboard.
\end{enumerate}

\begin{quote}
Field requirements and validation messages are shown inline if input is invalid.
\end{quote}

\subsection{Log In}
\begin{enumerate}
  \item Open CodeTrail and select \textbf{Log In}.
  \item Enter your registered email/username and password (or use Google sign-in if enabled).
  \item On success, you are redirected to the \textbf{Home Dashboard}.
\end{enumerate}

\textbf{Session behavior}
\begin{itemize}
  \item Your session token is stored locally to keep you signed in.
  \item If your session expires or becomes invalid, you may see a session/authentication error and be prompted to sign in again.
\end{itemize}

\subsection{Forgot Password \& Reset}
\begin{enumerate}
  \item On the login screen, select \textbf{Forgot Password}.
  \item Enter your registered email address and submit.
  \item If email service is configured, follow the reset link sent to your inbox.
  \item Set a new password on the reset page and log in again.
\end{enumerate}

\begin{quote}
If you do not receive an email, check spam/junk folders or contact your administrator (email delivery depends on deployment configuration).
\end{quote}

\subsection{Log Out}
Use the \textbf{Logout/Profile} controls available in the application (e.g., sidebar/profile menu depending on view). Logging out clears the active session and returns you to the public/landing experience.

% ==========================================
\section{Home Dashboard}

The Home Dashboard is your central hub after login. The browser tab title reflects your context (e.g., user and/or active workspace name).

\subsection{Top Banner (Recent/Active Workspace)}
\begin{itemize}
  \item Shows your most recent/active workspace (title, language, last edited time, description).
  \item \textbf{Jump Back In:} Opens the selected workspace in the Modular Workspace.
  \item If no recent workspace is available, the banner prompts you to create a workspace.
\end{itemize}

\subsection{Metrics Overview}
Dashboard cards summarize key counts (examples):
\begin{itemize}
  \item \textbf{Workspaces} (total)
  \item \textbf{Collaborators} (unique across your workspaces)
  \item \textbf{Project Files} (aggregated)
  \item \textbf{Total Lines Synced / Activity} (based on recorded telemetry/history)
\end{itemize}

\begin{quote}
Metrics reflect available data from your account/workspaces and may show loading or empty states.
\end{quote}

\subsection{Sidebar Navigation}
The collapsible/hover-expand sidebar provides quick access to:
\begin{itemize}
  \item \textbf{Dashboard:} Return to top/home view.
  \item \textbf{Workspaces:} Open the Workspaces Directory.
  \item \textbf{Starred:} Open Workspaces Directory with the \textbf{Starred} filter applied.
  \item \textbf{Shared with me:} Filter/navigate to shared workspaces.
  \item \textbf{Archived:} Filter/navigate to archived workspaces.
  \item \textbf{Telemetry/Activity:} Jump to activity/timeline views (e.g., runs filter) where available.
  \item \textbf{Profile/Settings:} Open your profile/settings area.
  \item \textbf{Personal Goals:} Quick goal tracking widget (see 4.6).
\end{itemize}

Badges/counts on sidebar items update dynamically based on your data.

\subsection{Personal Timeline}
The \textbf{Personal Timeline} displays recent activity/events (workspace updates, code runs, saves/edits, joins/sessions, tasks/goals, etc.) in chronological context.

\textbf{Controls}
\begin{itemize}
  \item \textbf{Search:} Filter events by description, file name, workspace, and related fields (debounced search).
  \item \textbf{Category filters:} e.g., All, Workspace, Runs/Code Execution, Tasks.
  \item \textbf{Workspace filter:} Limit events to a specific workspace.
  \item \textbf{Date range:} e.g., All Time, This Week, Today.
  \item \textbf{Reset/Clear filters:} Restore default view.
\end{itemize}

\textbf{Viewing behavior}
\begin{itemize}
  \item Shows \textbf{recent activity first} with a manageable initial set.
  \item Use \textbf{Show more} to load additional events; \textbf{Show less} collapses back to recent items.
  \item The header indicates total/visible event counts.
\end{itemize}

\textbf{Interaction}
\begin{itemize}
  \item Click/tap a timeline item to inspect it and (when applicable) jump to the related workspace.
\end{itemize}

\subsection{Recent Workspaces Panel}
\begin{itemize}
  \item Displays your recent workspaces with quick actions (e.g., open, copy room/invite code, view details).
  \item Filter pills (All/Active/Starred/Shared/Archived) help narrow the list.
  \item \textbf{New Workspace} and \textbf{Join} actions are available from this panel.
\end{itemize}

\subsection{Personal Goals}
\begin{itemize}
  \item Add, complete, or remove personal goals from the sidebar widget.
  \item Progress indicator updates based on completed goals.
  \item Goals are stored locally per user profile/session context in the browser.
\end{itemize}

% ==========================================
\section{Workspaces Directory}

The Workspaces Directory is the full management view for your workspaces.

\subsection{Open the Directory}
\begin{itemize}
  \item Click \textbf{Workspaces} in the sidebar or \textbf{View All} from the dashboard.
  \item Clicking \textbf{Starred} in the sidebar opens the directory with the \textbf{Starred} filter pre-selected.
\end{itemize}

\subsection{Search and Filters}
\begin{itemize}
  \item \textbf{Search:} Find workspaces by name, description, or room/invite code.
  \item \textbf{Category tabs/filters:} \textbf{All}, \textbf{Starred}, \textbf{Shared}, \textbf{Archived} (and other available categories).
  \item \textbf{Language filter:} Narrow by programming language (if available).
  \item \textbf{View options:} Grid/list view (where supported).
\end{itemize}

\subsection{Create a Workspace}
\begin{enumerate}
  \item Click \textbf{New Workspace} / \textbf{Create Workspace}.
  \item Provide required details (e.g., title, language, description/visibility/settings as available).
  \item Confirm to create. The workspace appears in your list and can be launched immediately.
\end{enumerate}

\subsection{Join a Workspace via Code}
\begin{enumerate}
  \item Click \textbf{Join via Code} / \textbf{Join}.
  \item Enter the workspace room/invite code shared by a collaborator.
  \item Submit to join (permissions depend on workspace settings and your role).
\end{enumerate}

\subsection{Star/Unstar a Workspace}
\begin{itemize}
  \item Click the \textbf{Star} icon on a workspace card/row to mark it as favorite.
  \item Starred state is saved for your account (persisted) and reflected across views (Home, Directory, sidebar counts).
  \item Use the \textbf{Starred} filter/tab to view only starred workspaces.
\end{itemize}

\subsection{Launch, Manage, and View Workspace Details}
\begin{itemize}
  \item \textbf{Launch IDE:} Open the workspace in the Modular Workspace.
  \item \textbf{Copy room/invite code:} Share access with collaborators.
  \item \textbf{Settings/management:} Open workspace settings (rename/configure/archive/permissions as available and permitted by your role).
  \item Workspace cards typically show language, privacy/status, file/member counts, and last activity indicators.
\end{itemize}

% ==========================================
\section{Modular Workspace (Code Editor)}

The Modular Workspace is the collaborative coding environment.

\subsection{Workspace Layout}
Common areas include:
\begin{itemize}
  \item \textbf{Editor:} Main code editing area (syntax highlighting, tabs/files).
  \item \textbf{File explorer/sidebar:} Navigate and manage project files/folders.
  \item \textbf{Top/bottom toolbars \& panels:} Run controls, status indicators, output/console, telemetry/audit tabs.
  \item \textbf{Collaboration/session info:} Room/workspace identifiers and participant context.
\end{itemize}

\begin{quote}
The browser tab displays the workspace name and user context when a workspace is open.
\end{quote}

\subsection{Files and Folders}
\begin{itemize}
  \item Open files from the explorer to edit.
  \item Create, rename, or delete files/folders (subject to permissions).
  \item File changes may auto-save or sync based on workspace/editor behavior and settings.
\end{itemize}

\subsection{Editing Code}
\begin{itemize}
  \item Use the editor for standard code editing (typing, selection, copy/paste).
  \item Language-aware behavior (highlighting/formatting cues) depends on file type/language.
  \item Unsaved/synced state indicators may appear based on configuration.
\end{itemize}

\subsection{Real-Time Collaboration}
\begin{itemize}
  \item Multiple users can collaborate in the same workspace session (when supported/enabled).
  \item Presence, edits, and activity may sync in real time via live connections.
  \item Connection/sync status indicators (e.g., live socket sync) may be shown in relevant panels.
\end{itemize}

\subsection{Running Code \& Execution Output}
\begin{enumerate}
  \item Open/select the file or code you want to run.
  \item Click \textbf{Run} (or the execution control provided in the workspace).
  \item The system sends code to the execution service and displays results in the output/console area.
\end{enumerate}

\textbf{Execution lifecycle states}
\begin{itemize}
  \item \textbf{Idle:} Ready before execution.
  \item \textbf{Running:} Execution in progress.
  \item \textbf{Success:} Completed with exit code \texttt{0}.
  \item \textbf{Error:} Runtime failure (non-zero exit code / stderr).
  \item \textbf{Compile Error:} Compilation failed (for compiled languages).
  \item \textbf{Timeout:} Execution exceeded allowed time limit.
\end{itemize}

\textbf{Result details}
\begin{itemize}
  \item \textbf{stdout:} Standard program output (formatted/monospace).
  \item \textbf{stderr/runtime errors:} Displayed separately (e.g., ``Runtime Error'').
  \item \textbf{Compilation errors:} Compiler diagnostics shown separately (e.g., ``Compilation Error'').
  \item \textbf{Exit code \& duration:} Reported when available (e.g., \texttt{Exit 0 • 24ms}).
  \item \textbf{Empty output:} Successful runs with no stdout show a clear ``no standard output produced'' notice.
\end{itemize}

\subsection{Output/Console Panel}
\begin{itemize}
  \item View execution results, logs, and status messages.
  \item Resize or toggle/minimize the panel (where supported) to optimize editor space.
  \item Some workspaces may include stdin/input support for programs that read input.
\end{itemize}

\subsection{Telemetry \& Audit Panel}
\begin{itemize}
  \item Displays contribution/activity metrics and integrity/audit information for the workspace/session.
  \item May include verification status, block/event summaries, and export/inspect actions depending on permissions and available data.
\end{itemize}

\subsection{AI Assistant (if enabled)}
If AI features are configured for your deployment:
\begin{itemize}
  \item Open the \textbf{AI Assistant} panel from the workspace.
  \item Use quick actions or chat to request explanations, refactoring, documentation, debugging help, tests, or code generation.
  \item Provide context (active file, selection, or attached files) as required by the assistant UI.
  \item Review suggested changes carefully; apply previews/diffs only after confirmation.
\end{itemize}

\begin{quote}
AI availability, models, and response quality depend on server configuration and provider status. When unavailable, the system may fall back to limited/local responses or show an error.
\end{quote}

% ==========================================
\section{Contribution Telemetry \& Proof-of-Work Audit}

The \textbf{Contribution Telemetry \& Integrity Auditor} provides visibility into recorded workspace activity and integrity verification.

\textbf{Typical elements}
\begin{itemize}
  \item \textbf{Status badges:} Audit/verification state and live sync indicators.
  \item \textbf{Summary metrics:} Contributors, lines synced, blocks/events, session/workspace identifiers.
  \item \textbf{Verification status:} Cryptographic/chain verification result (e.g., verified or discrepancy/tamper warning) when data is available.
  \item \textbf{Contributor table:} Responsibility/role, lines, participation, and verification/hash details.
  \item \textbf{Actions (as available):} Refresh/audit chain, inspect ledger, simulate/restore (demo/test flows), export reports.
\end{itemize}

\begin{quote}
Displayed data depends on workspace selection, permissions, and recorded telemetry.
\end{quote}

% ==========================================
\section{Contribution Analytics}

Analytics views provide higher-level insights into collaboration and contributions, such as:
\begin{itemize}
  \item Session/activity summaries and timelines
  \item Contributor participation and category breakdowns
  \item Execution/collaboration statistics
\end{itemize}

Use filters and workspace selectors (where provided) to focus on specific periods or contributors.

% ==========================================
\section{Profile \& Settings}

From the \textbf{Profile/Settings} area you can typically:
\begin{itemize}
  \item View and update profile information (as supported)
  \item Manage account/session preferences
  \item Access application-level settings exposed by your deployment
\end{itemize}

Available options depend on your role and enabled features.

% ==========================================
\section{Troubleshooting \& FAQs}

\subsection{I cannot log in / I get ``Session expired''}
\begin{itemize}
  \item Sign in again to refresh your session.
  \item Ensure cookies/local storage are enabled in your browser.
  \item If the issue persists, the backend authentication service or configuration may need attention.
\end{itemize}

\subsection{Workspaces do not load / ``Database Offline'' or connection error}
\begin{itemize}
  \item This indicates the app cannot reach required backend data services.
  \item Check your internet connection and retry.
  \item If running a local/dev environment, verify backend server and database are running and configured correctly.
\end{itemize}

\subsection{Code execution fails or never completes}
\begin{itemize}
  \item Confirm the execution service (sandbox/runtime) is available.
  \item Check for compile/runtime errors, timeouts, or unsupported language/runtime in your environment.
  \item Retry after service availability is restored.
\end{itemize}

\subsection{I don't see my starred workspace in Starred filter}
\begin{itemize}
  \item Ensure you starred the workspace while signed in (starring is account-specific and persisted).
  \item Refresh the directory; use \textbf{All} to confirm the workspace exists, then star it again if needed.
\end{itemize}

\subsection{Starred sidebar opens directory but filter looks wrong}
\begin{itemize}
  \item Use the \textbf{Starred} tab/filter in the Workspaces Directory to confirm.
  \item Refresh the page to reload latest workspace data.
\end{itemize}

\subsection{AI assistant is unavailable or gives generic/local responses}
\begin{itemize}
  \item AI features require valid provider configuration/credentials on the backend.
  \item If providers are unavailable or rate-limited, responses may fall back or fail; try again later or contact your administrator.
\end{itemize}

\subsection{Real-time collaboration/sync not updating}
\begin{itemize}
  \item Check network stability and live connection indicators.
  \item Refresh the workspace; if issues persist, verify backend/socket service availability.
\end{itemize}

% ==========================================
\section{Tips for Best Experience}
\begin{itemize}
  \item Keep your browser updated for best editor/performance support.
  \item Use descriptive workspace names and keep room/invite codes private when needed.
  \item Star frequently used workspaces for quick access.
  \item Save and run code regularly; review execution output (stdout/stderr/exit code) to diagnose issues.
  \item Use timeline filters and search to quickly locate recent work.
  \item For collaborative sessions, coordinate roles/permissions to avoid accidental edits.
\end{itemize}

% ==========================================
\section{Glossary}
\begin{description}
  \item[Workspace:] A collaborative coding environment/project room in CodeTrail.
  \item[Room/Invite Code:] A code used to join a workspace.
  \item[Starred Workspace:] A workspace marked as favorite for quick filtering/access.
  \item[stdout:] Standard output produced by a program.
  \item[stderr:] Error output produced during program execution.
  \item[Exit Code:] Numeric status returned by a process (\texttt{0} typically indicates success).
  \item[Timeout:] Execution stopped after exceeding the allowed runtime limit.
  \item[Telemetry:] Recorded activity/events used for contribution tracking and insights.
  \item[Proof-of-Work / Integrity Audit:] Verification view for recorded activity chain integrity (as implemented).
  \item[Session:] A period of user/workspace activity; may be tracked in timeline/history.
\end{description}

\begin{center}
\textit{End of CodeTrail User Manual}
\end{center}

\end{document}
